\scanpage{284}
\mspage{200}
% archive sheet number: 283
de cardinal $\leq\unsure{\pi|\kappa|k}_0$, représentables dans $M$, et notamment par fact. \unsure{directs}, \deleted{\unsure{grdes} devant $\pi_0$}, mais même alors, et $=$ en admettant que $M_{\pi_0}$ lui-$=$\unsure{!} est équivalent à une petite catégorie\deleted{\unsure{)}} , et prenant $C=M_{\pi_0}$\deleted{\unsure{)}} , il n'est pourtant pas clair que \deleted{\unsure{le f.}} \inserted{ce fonct. pl. fid.} $\mathrm{Ind}_{\pi_0} C\longrightarrow M$ soit ess. surjectif), ce qui équiv\deleted{\unsure{alvait|alait}} revient à dire que pour \unsure{tt|$\forall$} $\deleted{\unsure{E}}F$ dans $M$, $C_{/F}$ est \unsure{grde} devant $\pi_0$. Il n'est en effet pas $=$ clair que $C_{/\unsure{F|P}}$ soit filtrante.

Donc finalement, je doute qu'on ait a) $\Longrightarrow$ c') ou seulement a) $\Longrightarrow$ b), la situation n'est pas bien comprise visiblement.

Mais prouvons quand même

\[ b\Longrightarrow c'. \]

Pour ceci, \unsure{OPS|WPS} \deleted{\unsure{par}} $M=\mathrm{Ind}_{\pi_0}C$, $C$ petite. On sait alors que la filtration de $M$ par les $(M_{\unsure{\pi_0|\pi}})_{\pi\geq\pi_0}$ est une \deleted{\unsure{filtrati}} fega (prop. 4.7.7 et déf. \unsure{4.7.8'|4.7.8}), et $=$ une feg \uline{canonique}. De plus, je \inserted{dis que} le cardinal $\pi_0$ est \uline{utile} \inserted{et $=$ régulier} pour cette filtration. Ce qui achèvera de prouver c'). \deleted{En effet,} On a un résultat un peu plus précis\unsure{:}
